\documentclass[11pt,a4paper]{article}
\usepackage[italian]{babel}
\usepackage[margin=2.5cm]{geometry}
\usepackage{parskip}
\usepackage{amsmath, amssymb, amsthm}
\usepackage{booktabs}
\usepackage{array}
\usepackage{xcolor}
\usepackage{needspace}
\usepackage{enumitem}
\usepackage{listings}
\usepackage{tikz}
\usepackage[most]{tcolorbox}
\usepackage[hidelinks]{hyperref}
\usetikzlibrary{decorations.pathreplacing, shapes.geometric, arrows.meta, positioning}

% Niente vedove/orfane: i paragrafi non lasciano righe isolate a fine/inizio pagina
\widowpenalty=10000
\clubpenalty=10000
\linespread{0.99}

% Elenchi più compatti
\setlist{itemsep=2pt, parsep=0pt, topsep=4pt, partopsep=0pt}

% --- Riquadri con bordo nero, senza numero (si spezzano tra due pagine) ---
% Uso: \begin{definizione} ... \end{definizione}
%      \begin{teorema}[Nome] ... \end{teorema}  (il nome è facoltativo)
\tcbset{nota/.style={enhanced jigsaw, breakable, colback=white,
  colframe=black, boxrule=0.5pt, sharp corners}}
\NewTColorBox{definizione}{}{nota,
  before upper={\textbf{Definizione.}\ }}
\NewTColorBox{teorema}{o}{nota,
  before upper={\textbf{Teorema\IfValueT{#1}{ (#1)}.}\ }}
\NewTColorBox{proposizione}{o}{nota,
  before upper={\textbf{Proposizione\IfValueT{#1}{ (#1)}.}\ }}
\NewTColorBox{proprieta}{o}{nota,
  before upper={\textbf{Proprietà\IfValueT{#1}{ (#1)}.}\ }}
\NewTColorBox{assioma}{o}{nota, boxrule=1pt,
  before upper={\textbf{Assioma\IfValueT{#1}{ (#1)}.}\ }}

% --- Ambienti semplici (senza riquadro) ---
% Il titolo sta in cima al blocco, anche se il contenuto inizia con un riquadro o
% con codice; e non resta mai da solo in fondo a una pagina.
% Uso: \begin{esempio}[Titolo facoltativo] ... \end{esempio}
\newcommand{\newrunin}[2]{%
  \NewDocumentEnvironment{#1}{o}{%
    \par\Needspace{4\baselineskip}\addvspace{\medskipamount}\noindent
    \textbf{#2}\IfValueT{##1}{ (##1)}\textbf{.}\ \ignorespaces}%
  {\par\addvspace{\medskipamount}}}
\newrunin{esempio}{Esempio}
\newrunin{esercizio}{Esercizio}
\NewTColorBox{osservazione}{}{nota,
  before upper={\textbf{Osservazione.}\ }}

% V e F nelle tavole di verità
\newcommand{\V}{\textsf{V}}
\newcommand{\F}{\textsf{F}}

% --- Codice C ---
% Uso: \begin{codice} ... \end{codice}      codice C numerato
%      \begin{comando} ... \end{comando}    comandi da terminale
%      \begin{risultato} ... \end{risultato} output di un programma
%      \lstinline|printf("ciao");|           codice dentro al testo
\lstdefinestyle{c}{
  language=C,
  basicstyle=\ttfamily\small,
  keywordstyle=\color{blue}\bfseries,
  commentstyle=\color{gray}\itshape,
  stringstyle=\color{red!70!black},
  numbers=left,
  numberstyle=\tiny\color{gray},
  numbersep=8pt,
  columns=flexible,
  keepspaces=true,
  breaklines=true,
  showstringspaces=false,
  tabsize=4,
  morekeywords={include, define},
  % lettere accentate nei commenti
  literate={à}{{\`a}}1 {è}{{\`e}}1 {é}{{\'e}}1 {ì}{{\`i}}1
           {ò}{{\`o}}1 {ù}{{\`u}}1 {È}{{\`E}}1 {À}{{\`A}}1
}
\lstset{style=c}
\newtcblisting{codice}{listing only, nota, left=6mm,
  listing options={style=c}}
\newtcblisting{comando}{listing only, nota,
  listing options={basicstyle=\ttfamily\small, numbers=none, language={},
    columns=flexible, keepspaces=true, breaklines=true}}
\newtcblisting{risultato}{listing only, enhanced jigsaw, breakable,
  colback=black!5, colframe=black, boxrule=0.5pt, sharp corners,
  listing options={basicstyle=\ttfamily\small, numbers=none, language={},
    columns=flexible, keepspaces=true, breaklines=true}}

% --- Diagrammi di flusso (simboli standard) ---
\tikzset{
  terminale/.style = {ellipse, draw, minimum width=2.5cm, minimum height=0.9cm},
  io/.style        = {trapezium, trapezium left angle=70, trapezium right angle=110,
                      draw, minimum width=2.5cm, minimum height=0.9cm},
  processo/.style  = {rectangle, draw, minimum width=2.5cm, minimum height=0.9cm},
  decisione/.style = {diamond, draw, aspect=2, minimum width=2.5cm},
  freccia/.style   = {thick, -{Stealth}}
}

\title{Programmazione procedurale\\[0.4em]\large Appunti}
\author{Marco Cerbella}
\date{Università degli Studi di Perugia -- A.A. 2026/2027\\ Aggiornato al \today}

\begin{document}
\maketitle
\setcounter{tocdepth}{1}
\tableofcontents
\clearpage

\section[Algoritmi e diagrammi di flusso (Lezione 1 -- 22 settembre 2026)]{Algoritmi e diagrammi di flusso}
\noindent\emph{Lezione 1 -- 22 settembre 2026}
\subsection{Algoritmi e diagrammi di flusso}

\begin{definizione}
Un \emph{algoritmo} è una procedura che risolve un problema attraverso una
sequenza finita di passi elementari e non ambigui. Spesso comporta la
ripetizione di alcune operazioni.
\end{definizione}

\begin{esempio}
Il primo algoritmo celebre è quello di Euclide (\emph{Elementi}, 300 a.C.) per
il massimo comun divisore.
\end{esempio}

\begin{definizione}
Un \emph{diagramma di flusso} (\emph{flowchart}) è la rappresentazione grafica
di un algoritmo. Le operazioni si leggono dall'alto verso il basso, seguendo le
frecce. I simboli sono standardizzati (ANSI/ISO).
\end{definizione}

I simboli principali sono:
\begin{itemize}
    \item \textbf{ellisse}: inizio e fine;
    \item \textbf{parallelogramma}: input e output (I/O);
    \item \textbf{rettangolo}: operazione (\emph{processo});
    \item \textbf{rombo}: decisione (condizione vero/falso).
\end{itemize}

\begin{center}
\begin{tikzpicture}[node distance=0.8cm]
  \node (inizio) [terminale] {Inizio};
  \node (input)  [io, below=of inizio] {Input};
  \node (op)     [processo, below=of input] {Operazioni};
  \node (output) [io, below=of op] {Output};
  \node (fine)   [terminale, below=of output] {Fine};

  \draw [freccia] (inizio) -- (input);
  \draw [freccia] (input) -- (op);
  \draw [freccia] (op) -- (output);
  \draw [freccia] (output) -- (fine);
\end{tikzpicture}
\end{center}

Lo stesso schema in linguaggio \textit{C}:
\begin{codice}
// Importare le librerie necessarie
#include <stdio.h>

// Funzione principale
int main() {
    // Lettura dell'input e operazioni

    // Stampa del risultato
    return 0;
}
\end{codice}

\subsection{Programmare}

\begin{definizione}
\emph{Programmare} significa tradurre un algoritmo in una notazione, un
linguaggio di programmazione, che un elaboratore può eseguire. Un
\emph{programma} è un insieme di istruzioni.
\end{definizione}

I passi sono: partire dal problema, \emph{pensare all'algoritmo}, scegliere il
linguaggio giusto e solo dopo programmare.

\subsection{Esempio: massimo comun divisore}

Algoritmo di Euclide (versione con sottrazioni): finché $b \neq 0$, si
sostituisce il maggiore tra $a$ e $b$ con la differenza tra i due. Alla fine $a$
contiene il MCD.

\begin{center}
\begin{tikzpicture}
  \node (s) at (0,0) [terminale] {Inizio};
  \node (in) at (0,-1.3) [io] {Leggi $a$, $b$};
  \node (d1) at (0,-3.3) [decisione] {$b \neq 0$};
  \node (d2) at (0,-5.8) [decisione] {$a > b$};
  \node (pa) at (-2.3,-8) [processo] {$a = a - b$};
  \node (pb) at (2.3,-8) [processo] {$b = b - a$};
  \node (j) at (0,-9.5) {};
  \node (m) at (0,-2.3) {};
  \node (out) at (4.3,-3.3) [io] {Stampa $a$};
  \node (e) at (4.3,-4.8) [terminale] {Fine};

  \draw [freccia] (s) -- (in);
  \draw [freccia] (in) -- (d1);
  \draw [freccia] (d1) -- node[right]{Sì} (d2);
  \draw [freccia] (d1.east) -- node[above]{No} (out.west);
  \draw [freccia] (out) -- (e);
  \draw [freccia] (d2.west) -| node[above, pos=0.2]{Sì} (pa.north);
  \draw [freccia] (d2.east) -| node[above, pos=0.2]{No} (pb.north);
  \draw (pa.south) |- (j.center);
  \draw (pb.south) |- (j.center);
  \draw [freccia] (j.center) -- (-4.4,-9.5) -- (-4.4,-2.3) -- (m.center);
\end{tikzpicture}
\end{center}

E il programma C corrispondente:
\begin{codice}
#include <stdio.h>

int main() {
    int a, b;

    printf("Enter first positive integer: \n");
    scanf("%d", &a);
    printf("Enter second positive integer: \n");
    scanf("%d", &b);

    while (b != 0) {
        if (a > b)
            a = a - b;
        else
            b = b - a;
    }
    printf("GCD = %d\n", a);

    return 0;
}
\end{codice}

\begin{osservazione}
Il ciclo \lstinline|while| corrisponde al rombo $b \neq 0$ con la freccia che
torna indietro; l'\lstinline|if|/\lstinline|else| corrisponde al secondo rombo.
\end{osservazione}

\section[Esercizi sui diagrammi di flusso (Lezione 2 -- 23 settembre 2026)]{Esercizi sui diagrammi di flusso}
\noindent\emph{Lezione 2 -- 23 settembre 2026}
\subsection{Somma di 10 e 20}

Si usa una variabile \texttt{sum} inizializzata a zero; si leggono i due numeri,
si sommano salvando il risultato in \texttt{sum} e infine si stampa.

\begin{enumerate}
    \item Inizializza $sum = 0$ (processo)
    \item Leggi i numeri (I/O)
    \item $sum = 10 + 20$ (processo)
    \item Stampa $sum$ (I/O)
\end{enumerate}

\begin{center}
\begin{tikzpicture}[node distance=0.7cm]
  \node (a) [terminale] {Inizio};
  \node (b) [processo, below=of a] {$sum = 0$};
  \node (c) [io, below=of b] {Leggi i numeri};
  \node (d) [processo, below=of c] {$sum = 10 + 20$};
  \node (e) [io, below=of d] {Stampa $sum$};
  \node (f) [terminale, below=of e] {Fine};
  \draw [freccia] (a) -- (b);
  \draw [freccia] (b) -- (c);
  \draw [freccia] (c) -- (d);
  \draw [freccia] (d) -- (e);
  \draw [freccia] (e) -- (f);
\end{tikzpicture}
\end{center}

\subsection{Somma di 5 numeri}

Servono due variabili: \texttt{sum} (il risultato) e \texttt{count} (quanti
numeri sono stati letti). Si usa un \emph{ciclo}: i passi si ripetono finché la
condizione è vera.

\begin{enumerate}
    \item $sum = 0$, $count = 0$ (processo)
    \item Leggi $n$ (I/O)
    \item $sum = sum + n$ e $count = count + 1$ (processo)
    \item $count < 5$? (decisione)
    \item Se sì, torna al passo 2; altrimenti stampa $sum$ (I/O)
\end{enumerate}

\begin{center}
\begin{tikzpicture}
  \node (a) at (0,0) [terminale] {Inizio};
  \node (b) at (0,-1.4) [processo] {$sum = 0,\ count = 0$};
  \node (c) at (0,-2.8) [io] {Leggi $n$};
  \node (d) at (0,-4.2) [processo] {$sum = sum + n$};
  \node (d2) at (0,-5.4) [processo] {$count = count + 1$};
  \node (e) at (0,-7.2) [decisione] {$count < 5$};
  \node (f) at (0,-9.2) [io] {Stampa $sum$};
  \node (g) at (0,-10.6) [terminale] {Fine};
  \draw [freccia] (a) -- (b);
  \draw [freccia] (b) -- (c);
  \draw [freccia] (c) -- (d);
  \draw [freccia] (d) -- (d2);
  \draw [freccia] (d2) -- (e);
  \draw [freccia] (e) -- node[right]{No} (f);
  \draw [freccia] (f) -- (g);
  \draw [freccia] (e.west) -- node[above]{Sì} ++(-2.3,0) |- (c.west);
\end{tikzpicture}
\end{center}

Il programma C corrispondente:
\begin{codice}
#include <stdio.h>

int main() {
    int sum = 0, count = 0, n;
    while (count < 5) {
        scanf("%d", &n);
        sum = sum + n;
        count = count + 1;
    }
    printf("sum = %d\n", sum);
    return 0;
}
\end{codice}

\subsection{Stampare ``Hello World'' 10 volte}

Anche qui si usa un ciclo con un contatore.

\begin{enumerate}
    \item $count = 0$ (processo)
    \item Stampa ``Hello World'' (I/O)
    \item $count = count + 1$ (processo)
    \item $count < 10$? Se sì torna al passo 2, altrimenti termina.
\end{enumerate}

\begin{center}
\begin{tikzpicture}
  \node (a) at (0,0) [terminale] {Inizio};
  \node (b) at (0,-1.4) [processo] {$count = 0$};
  \node (c) at (0,-2.8) [io] {Stampa ``Hello World''};
  \node (d) at (0,-4.2) [processo] {$count = count + 1$};
  \node (e) at (0,-6) [decisione] {$count < 10$};
  \node (f) at (0,-8) [terminale] {Fine};
  \draw [freccia] (a) -- (b);
  \draw [freccia] (b) -- (c);
  \draw [freccia] (c) -- (d);
  \draw [freccia] (d) -- (e);
  \draw [freccia] (e) -- node[right]{No} (f);
  \draw [freccia] (e.west) -- node[above]{Sì} ++(-2.6,0) |- (c.west);
\end{tikzpicture}
\end{center}

\begin{codice}
#include <stdio.h>

int main() {
    int count = 0;
    do {
        printf("Hello World\n");
        count = count + 1;
    } while (count < 10);
    return 0;
}
\end{codice}

\begin{osservazione}
Il diagramma esegue la stampa \emph{prima} di controllare la condizione, quindi
il costrutto C fedele è \lstinline|do ... while|. Con un \lstinline|while|
la condizione verrebbe controllata prima del corpo.
\end{osservazione}

\subsection{Login a un account (da risolvere liberamente)}

Esercizio aperto: non esiste una sola soluzione. Una possibile sequenza:
\begin{enumerate}
    \item Inserisci l'indirizzo del sito nel browser (I/O)
    \item Si carica la home page (processo)
    \item Inserisci email e password (I/O)
    \item Credenziali valide? (decisione)
    \item Se no: errore di login (processo) e torna al passo 3;
          se sì: mostra l'account (I/O) e termina.
\end{enumerate}

\begin{center}
\begin{tikzpicture}
  \node (a) at (0,0) [terminale] {Inizio};
  \node (b) at (0,-1.3) [io] {Inserisci URL};
  \node (c) at (0,-2.6) [processo] {Si carica la home};
  \node (d) at (0,-3.9) [io] {Inserisci email e password};
  \node (e) at (0,-5.8) [decisione] {Credenziali valide?};
  \node (f) at (0,-7.8) [io] {Mostra l'account};
  \node (g) at (0,-9.1) [terminale] {Fine};
  \node (err) at (4.8,-5.8) [processo] {Errore di login};
  \draw [freccia] (a) -- (b);
  \draw [freccia] (b) -- (c);
  \draw [freccia] (c) -- (d);
  \draw [freccia] (d) -- (e);
  \draw [freccia] (e) -- node[right]{Sì} (f);
  \draw [freccia] (f) -- (g);
  \draw [freccia] (e.east) -- node[above]{No} (err.west);
  \draw [freccia] (err.north) |- (d.east);
\end{tikzpicture}
\end{center}

\section[Il linguaggio C (Lezione 3 -- 28 settembre 2026)]{Il linguaggio C}
\noindent\emph{Lezione 3 -- 28 settembre 2026}
\subsection{Linguaggi di programmazione}

\begin{definizione}
Un \emph{linguaggio di programmazione} è un linguaggio formale che specifica un
insieme di istruzioni per un elaboratore. Si usa per realizzare programmi che
implementano algoritmi. I linguaggi \emph{ad alto livello} sono vicini al
linguaggio umano e lontani dal linguaggio macchina.
\end{definizione}

I \emph{paradigmi di programmazione} classificano i linguaggi in base alle loro
caratteristiche. Il C è \emph{imperativo}: il programmatore indica alla
macchina \emph{come} modificare il proprio stato, ad esempio
\lstinline|a = a + 1;|.

\subsection{Dal linguaggio macchina al C}

\begin{itemize}
    \item Il \emph{linguaggio macchina} (sequenze di bit) è l'unico eseguito
          direttamente dal processore.
    \item L'\emph{assembly} usa abbreviazioni mnemoniche per le operazioni,
          ma è legato alla macchina.
    \item I linguaggi \emph{ad alto livello} sono indipendenti dalla macchina.
\end{itemize}

\begin{definizione}
Un \emph{compilatore} è un programma di sistema che traduce il programma
sorgente scritto in un linguaggio ad alto livello in un programma
\emph{target} equivalente (assembly/linguaggio macchina). Il programma target
viene eseguito in un secondo momento.
\end{definizione}

In questo corso il compilatore è \textbf{gcc}.

\begin{center}
\begin{tikzpicture}[node distance=1.6cm]
  \node (src) [processo] {sorgente C};
  \node (gcc) [processo, right=of src] {gcc};
  \node (exe) [processo, right=of gcc] {programma eseguibile};
  \draw [freccia] (src) -- (gcc);
  \draw [freccia] (gcc) -- (exe);
\end{tikzpicture}
\end{center}

\subsection{Caratteristiche del C}

Il C è un linguaggio ad alto livello, \emph{procedurale} e di uso generale,
progettato (Dennis Ritchie, anni '70) per scrivere sistemi operativi con la
massima indipendenza dall'hardware.

\textbf{Pregi:}
\begin{itemize}
    \item veloce (compilato, vicino alla macchina) e portabile;
    \item linguaggio piccolo e maturo, con molti strumenti;
    \item accesso diretto alla memoria e alle funzionalità di basso livello.
\end{itemize}

\textbf{Difficoltà:}
\begin{itemize}
    \item si può facilmente fare danni con accessi diretti alla memoria e
          puntatori;
    \item la memoria va gestita a mano;
    \item il codice è più verboso rispetto a linguaggi di scripting come Python.
\end{itemize}

\begin{osservazione}
Il linguaggio C in sé non ha istruzioni di input/output né di gestione della
memoria dinamica: queste funzioni sono fornite dalla \emph{libreria standard}
(ad esempio \lstinline|stdio.h|).
\end{osservazione}

\subsection{Standard}

Il C è stato standardizzato più volte: C89/C90 (ANSI/ISO), C99, C11, C18. Con
gcc si può scegliere lo standard:
\begin{comando}
gcc file.c -std=c11
\end{comando}

\subsection{Compilati e interpretati}

Un linguaggio \emph{compilato} (come C) viene analizzato a fondo dal traduttore
prima dell'esecuzione, e quasi tutti i controlli semantici sono fatti a tempo di
compilazione. In un linguaggio \emph{interpretato} (come Python) l'interprete
resta presente durante l'esecuzione e la maggior parte dei controlli è fatta a
tempo di esecuzione.

\section[Struttura di un programma C e compilazione (Lezione 4 -- 29 settembre 2026)]{Struttura di un programma C e compilazione}
\noindent\emph{Lezione 4 -- 29 settembre 2026}
\subsection{Struttura di un programma}

Gli elementi costitutivi di un programma C sono le \emph{funzioni}, che possono
chiamarsi a vicenda.
\begin{itemize}
    \item Ogni funzione ha uno scopo preciso; le funzioni \emph{non} si possono
          annidare.
    \item Una funzione contiene istruzioni eseguite in sequenza; le istruzioni
          si raggruppano in \emph{blocchi}.
    \item Si possono usare funzioni della libreria standard (\lstinline|printf|)
          o scriverne di proprie.
    \item Ogni programma deve definire almeno la funzione \lstinline|main()|:
          è la prima funzione eseguita ed è il livello principale di controllo.
\end{itemize}

\subsection{Hello World e compilazione con gcc}

\begin{codice}
#include <stdio.h>

int main() {
    printf("Hello World\n");
    return 0;
}
\end{codice}

Salvato il codice in \texttt{test.c}, da terminale si compila (creando l'eseguibile \texttt{test}) e poi si esegue:
\begin{comando}
gcc -o test test.c
./test
\end{comando}

\begin{osservazione}
Per il corso e per l'esame si compila \emph{solo con gcc} da terminale (anche
se si usa un editor che ha un terminale integrato). Senza \lstinline|-o| il
file eseguibile si chiama \texttt{a.out}.
\end{osservazione}

\subsection{Input e output: printf e scanf}

\begin{esempio}[Somma di due interi]
\begin{codice}
#include <stdio.h>
int main() {
    int firstNumber, secondNumber, sumOfTwoNumbers;

    printf("Enter two integers: ");

    // gli interi inseriti vengono memorizzati con scanf()
    scanf("%d %d", &firstNumber, &secondNumber);

    // la somma viene salvata in sumOfTwoNumbers
    sumOfTwoNumbers = firstNumber + secondNumber;

    // stampa della somma
    printf("%d + %d = %d", firstNumber, secondNumber, sumOfTwoNumbers);

    return 0;
}
\end{codice}
\begin{risultato}
Enter two integers: 12
11
12 + 11 = 23
\end{risultato}
\end{esempio}

\begin{esempio}[Quoziente e resto]
\begin{codice}
#include <stdio.h>
int main() {
    int dividend, divisor, quotient, remainder;

    printf("Enter dividend: ");
    scanf("%d", &dividend);

    printf("Enter divisor: ");
    scanf("%d", &divisor);

    quotient = dividend / divisor;    // quoziente
    remainder = dividend % divisor;   // resto

    printf("Quotient = %d\n", quotient);
    printf("Remainder = %d", remainder);

    return 0;
}
\end{codice}
\begin{risultato}
Enter dividend: 25
Enter divisor: 4
Quotient = 6
Remainder = 1
\end{risultato}
\end{esempio}

\begin{osservazione}
\lstinline|%d| è il segnaposto per un intero; in \lstinline|scanf| le variabili
vanno precedute da \lstinline|&|.
\end{osservazione}

\subsection{Programmazione procedurale e modulare}

\begin{definizione}
La \emph{programmazione procedurale} è un paradigma, derivato dalla
programmazione strutturata, basato sul concetto di chiamata di procedura. Le
procedure (o routine, sottoprogrammi, funzioni) contengono una serie di passi di
calcolo da eseguire. Il C è imperativo e procedurale: le istruzioni sono
raggruppate in funzioni.
\end{definizione}

\begin{esempio}
\texttt{mainfile.c}
\begin{codice}
#include <stdio.h>
void myPrintHello();

int main() {
    myPrintHello();
    return(0);
}

void myPrintHello(void) {
    printf("Hello!\n");
    return;
}
\end{codice}
\end{esempio}

\begin{definizione}
La \emph{programmazione modulare} è una tecnica di progettazione che separa le
funzionalità di un programma in \emph{moduli} indipendenti e intercambiabili,
ciascuno contenente tutto il necessario per realizzare un solo aspetto della
funzionalità.
\end{definizione}

Ad esempio il codice si può dividere in file diversi.

\texttt{mainfile.c}:
\begin{codice}
int main() {
    myPrintHello();
    return(0);
}
\end{codice}

\texttt{hello.c}:
\begin{codice}
#include <stdio.h>

void myPrintHello(void) {
    printf("Hello!\n");
    return;
}
\end{codice}

e si compilano insieme:
\begin{comando}
gcc -o executable mainfile.c hello.c
\end{comando}

\section[Errori di compilazione e di esecuzione (Lezione 5 -- 30 settembre 2026)]{Errori di compilazione e di esecuzione}
\noindent\emph{Lezione 5 -- 30 settembre 2026}
\subsection{Compile-time e run-time}

La vita di un programmatore ha due momenti: il \emph{compile-time}, quando si
esegue \texttt{gcc}, e il \emph{run-time}, quando si esegue il programma.

\begin{definizione}
Un \emph{errore di compilazione} (\emph{compile-time}) è un errore rilevato dal
compilatore. Cause comuni: errori di sintassi, come un punto e virgola mancante
o l'uso di una parola riservata (ad esempio \lstinline|while|) come nome.
\end{definizione}

\begin{definizione}
Un \emph{errore di esecuzione} (\emph{run-time}) è un errore del programma che
si verifica mentre il programma è in esecuzione.
\end{definizione}

\subsection{Esempio di errore di compilazione}

Nel programma del MCD manca il \texttt{;} dopo la seconda \lstinline|scanf|:
\begin{codice}
printf("Enter first positive integer: \n");
scanf("%d", &b)        // manca il ';'
\end{codice}

Il compilatore segnala:
\begin{risultato}
$ gcc -o euclid euclid.c
euclid.c:9:20: error: expected ';' after expression
   scanf("%d", &b)
                 ^
                 ;
1 error generated.
\end{risultato}

Il messaggio indica file, riga e colonna dell'errore: bisogna imparare a
leggere l'output del compilatore per eliminare tutti gli errori di compilazione.

\subsection{Esempio di errore di esecuzione}

Il programma compila senza problemi, ma dereferenzia un puntatore nullo:
\begin{codice}
#include <stdio.h>
int main() {
    int *a = NULL;

    printf("%d", *a);    // accesso a memoria non valida

    int i = 0;
    while (i < 4) {
        printf("Hello World");
        i++;
    }
    return 0;
}
\end{codice}
\begin{risultato}
$ gcc -o runtime_error runtime_error.c
$ ./runtime_error
Segmentation fault: 11
\end{risultato}

\begin{osservazione}
Eliminare gli errori di compilazione non basta: bisogna testare il programma
con molti input diversi per eliminare (il più possibile) gli errori di
esecuzione.
\end{osservazione}

\begin{center}
\emph{``Testing shows the presence, not the absence of bugs''} -- E.\,W.\ Dijkstra
\end{center}

\section[Commenti, identificatori e parole chiave (Lezione 6 -- 5 ottobre 2026)]{Commenti, identificatori e parole chiave}
\noindent\emph{Lezione 6 -- 5 ottobre 2026}
\subsection{Commenti}

In C ci sono due tipi di commenti:
\begin{itemize}
    \item \emph{a blocco}: iniziano con \lstinline|/*| e finiscono con \lstinline|*/|, possono
          occupare più righe;
    \item \emph{di riga}: iniziano con \lstinline|//| e finiscono a fine riga.
\end{itemize}

\begin{codice}
double pi = 3.1415926536;   // Pi è 3,14
\end{codice}

All'interno di una stringa o di una costante carattere i simboli \lstinline|/*| e
\lstinline|//| non iniziano un commento:
\begin{codice}
printf("Comments in C begin with /* or //.\n");
\end{codice}

\begin{osservazione}
I commenti a blocco \emph{non si possono annidare}: il primo \lstinline|*/|
chiude il commento aperto. Si possono invece usare per ``spegnere'' del codice
che contiene commenti di riga.
\end{osservazione}

\begin{esempio}
Qui il commento aperto alla riga 3 termina al primo \lstinline|*/| (riga 6): il
secondo \lstinline|*/| è un errore di sintassi.
\begin{codice}
#include <stdio.h>
int main() {
    int x = 4; /*
    int y = 10; /*

    printf("x = %d\n", x); */
    printf("y = %d\n", y); */
}
\end{codice}
\end{esempio}

\subsubsection{Scrivere buoni commenti}

Il codice si scrive per essere eseguito da un computer ma letto da persone:
da noi stessi mentre lo scriviamo e dopo settimane o mesi, dai colleghi che
devono integrarlo e da chi lo manterrà anni dopo.
\begin{itemize}
    \item Evitare commenti ridondanti, come \lstinline|i = i + 1; // incrementa i|.
    \item Non lasciare codice ``eliminato'' sotto forma di commento (va bene
          solo temporaneamente durante il debug).
    \item Evitare commenti decorati (riquadri di asterischi): sono difficili da
          mantenere.
    \item Un commento che chiarisce spesso segnala codice troppo complesso:
          meglio semplificare il codice, che deve essere \emph{autoesplicativo}.
\end{itemize}

\subsection{Caratteri}

Il C definisce due insiemi di caratteri: quello \emph{sorgente} (usabile nel
codice) e quello di \emph{esecuzione} (interpretabile dal programma in
esecuzione); nella maggior parte delle implementazioni coincidono. I caratteri
\emph{spazio bianco} sono: spazio, \verb|\f| (form feed), \verb|\n| (a capo),
\verb|\t| (tab orizzontale), \verb|\v| (tab verticale).

\subsection{Identificatori}

\begin{definizione}
Un \emph{identificatore} è il nome di una variabile, funzione, macro,
struttura o altro oggetto definito in un programma C. Può contenere:
\begin{itemize}
    \item lettere \texttt{a--z} e \texttt{A--Z} (il C è \emph{case-sensitive});
    \item il carattere underscore \texttt{\_};
    \item cifre \texttt{0--9}, ma il primo carattere \emph{non} può essere una cifra.
\end{itemize}
\end{definizione}

\begin{esempio}
Identificatori validi: \texttt{x}, \texttt{dollar}, \texttt{While},
\texttt{error\_handler}, \texttt{scale64}.

Non validi: \texttt{1st\_rank} (inizia con una cifra), \texttt{switch} (parola
riservata), \texttt{y/n} e \texttt{x-ray} (caratteri non ammessi).
\end{esempio}

Non c'è un limite alla lunghezza, ma i compilatori considerano significativi
solo i primi 31 o 63 caratteri. Molti nomi sono già usati dalla libreria
standard (ad esempio \lstinline|printf|) e non si possono riutilizzare per
funzioni proprie o variabili globali.

\subsection{Parole riservate}

Le seguenti 44 parole chiave sono riservate in C (erano 32 in C89), hanno un
significato preciso per il compilatore e non si possono usare come
identificatori:

\begin{center}
\begin{minipage}{0.9\textwidth}
\ttfamily\small
auto, break, case, char, const, continue, default, do, double, else, enum,
extern, float, for, goto, if, inline, int, long, register, restrict, return,
short, signed, sizeof, static, struct, switch, typedef, union, unsigned, void,
volatile, while, \_Alignas, \_Alignof, \_Atomic, \_Bool, \_Complex, \_Generic,
\_Imaginary, \_Noreturn, \_Static\_assert, \_Thread\_local
\end{minipage}
\end{center}

\subsubsection{Scegliere i nomi}

Un nome descrive l'oggetto, il suo comportamento e il suo uso: una variabile
chiamata male può confondere molto.
\begin{itemize}
    \item evitare nomi inutilmente lunghi;
    \item usare nomi diversi in ambiti diversi, per evitare confusione.
\end{itemize}

\begin{esempio}
Confrontare \lstinline|a = b * c;| con
\lstinline|weekly_pay = hours_worked * pay_rate;|.
\end{esempio}

\section[Scope e binding (Lezione 7 -- 6 ottobre 2026)]{Scope e binding}
\noindent\emph{Lezione 7 -- 6 ottobre 2026}
\subsection{Scope degli identificatori}

\begin{definizione}
Lo \emph{scope} (ambito) di un identificatore è la parte del programma in cui
l'identificatore è significativo, cioè in cui può essere ``visto''. Il tipo di
scope dipende da dove l'identificatore viene dichiarato.
\end{definizione}

\begin{itemize}
    \item \textbf{File scope}: l'identificatore è dichiarato fuori da tutti i
          blocchi e dalle liste di parametri (cioè fuori dalle funzioni). Si
          può usare dalla dichiarazione fino alla fine del file.
    \item \textbf{Block scope}: l'identificatore è dichiarato dentro un blocco.
          Si può usare dalla dichiarazione fino alla fine del blocco più piccolo
          che la contiene (spesso, ma non sempre, il corpo di una funzione).
\end{itemize}

\begin{esempio}
\begin{codice}
int p;                   // file scope

int fun2(void) {         // fun2: file scope
    int r = 3;           // r, s: block scope
    int s = r;
    while (r != s) {
        r--;
    }
    return r;
}
\end{codice}
\end{esempio}

\begin{osservazione}
Non si possono avere due variabili con lo stesso nome nello stesso scope:
\begin{codice}
int main() {
    int a = 5;
    float a = 6.5;   // errore: a già dichiarata
}
\end{codice}
\end{osservazione}

\subsection{Shadowing}

Si può riusare un identificatore in una dichiarazione \emph{annidata} nello
scope esistente; la nuova dichiarazione deve avere block scope, in un blocco
contenuto in quello esterno. In questo caso la dichiarazione interna
\emph{nasconde} (\emph{shadowing}) quella esterna, che non è visibile nello
scope interno.

\begin{esempio}
\begin{codice}
int main() {
    int x = 1, y = 2, z = 3;
    printf(" x = %d, y = %d, z = %d \n", x, y, z);
    {
        int x = 10;
        float y = 20;
        printf(" x = %d, y = %f, z = %d \n", x, y, z);
        {
            int z = 100;
            printf(" x = %d, y = %f, z = %d \n", x, y, z);
        }
    }
    return 0;
}
\end{codice}
\begin{risultato}
x = 1, y = 2, z = 3
x = 10, y = 20.000000, z = 3
x = 10, y = 20.000000, z = 100
\end{risultato}
\end{esempio}

\begin{esempio}
Qui \lstinline|x++| modifica la \lstinline|x| esterna (non è ridichiarata),
mentre \lstinline|y++| modifica la \lstinline|y| interna; uscendo dal blocco
interno la \lstinline|y| torna a valere 20.
\begin{codice}
int main() {
    {
        int x = 10, y = 20;
        {
            printf("x = %d, y = %d\n", x, y);
            {
                int y = 40;
                x++;
                y++;
                printf("x = %d, y = %d\n", x, y);
            }
            printf("x = %d, y = %d\n", x, y);
        }
    }
    return 0;
}
\end{codice}
\begin{risultato}
x = 10, y = 20
x = 11, y = 41
x = 11, y = 20
\end{risultato}
\end{esempio}

\begin{esempio}
Fuori dal blocco in cui è dichiarata, una variabile non è accessibile:
\begin{codice}
int main() {
    {
        int x = 10;
    }
    {
        printf("%d", x);    // errore: x non visibile qui
    }
    return 0;
}
\end{codice}
\begin{risultato}
error: 'x' undeclared (first use in this function)
\end{risultato}
\end{esempio}

\subsection{Binding statico}

\begin{definizione}
Un \emph{binding} è l'associazione tra un nome e l'oggetto che esso denota. Il
\emph{binding time} è il momento in cui il binding viene creato (più in
generale, il momento in cui viene presa una decisione di implementazione). La
regione testuale del programma in cui un binding è attivo è il suo scope.
\end{definizione}

In C lo scope di un binding è determinato \emph{staticamente}, cioè a tempo di
compilazione (\emph{static scoping}), come si vede dal testo del programma e non
dall'ordine delle chiamate a run-time. Altri linguaggi usano lo scoping
\emph{dinamico}: ad esempio in Perl \lstinline|my| definisce una variabile
locale a scope statico, \lstinline|local| una a scope dinamico.

\subsection{Esercizio}

\begin{esercizio}
Qual è l'output del seguente programma?
\begin{codice}
#include <stdio.h>

int f1(void);
int f2(int x, int a);

int a;

int main() {
    int a, b, c;

    a = 7;
    b = f1();
    c = f2(a, b);
    printf("%d %d %d\n", a, b, c);
}

int f1(void) {
    a = 12;
    printf("%d ", a);
    return (a + 5);
}

int f2(int x, int a) {
    printf("%d ", a);
    return (x * a);
}
\end{codice}
\end{esercizio}

\emph{Soluzione.}
\begin{itemize}
    \item In \lstinline|main| la \lstinline|a| locale (che vale 7) nasconde la
          \lstinline|a| globale.
    \item \lstinline|f1| non dichiara una propria \lstinline|a|: per binding
          statico modifica quella \emph{globale}, la pone a 12, stampa
          \texttt{12} e restituisce 17, quindi $b = 17$.
    \item In \lstinline|f2(7, 17)| il parametro \lstinline|a| vale 17: stampa
          \texttt{17} e restituisce $7 \cdot 17 = 119$, quindi $c = 119$.
    \item Infine \lstinline|main| stampa \texttt{7 17 119}.
\end{itemize}
\begin{risultato}
12 17 7 17 119
\end{risultato}

\section[Tipi di dato (Lezione 8 -- 7 ottobre 2026)]{Tipi di dato}
\noindent\emph{Lezione 8 -- 7 ottobre 2026}
\subsection{Tipi e oggetti}

I programmi devono memorizzare ed elaborare dati di tipo diverso, come interi e
numeri in virgola mobile, in modi diversi: il compilatore deve quindi sapere che
tipo di dato rappresenta un valore.

\begin{definizione}
In C un \emph{oggetto} è una zona di memoria il cui contenuto può rappresentare
valori. Gli oggetti che hanno un nome sono detti \emph{variabili}. Il
\emph{tipo} di un oggetto determina:
\begin{itemize}
    \item quanto spazio occupa in memoria;
    \item i valori che la variabile può assumere;
    \item le operazioni che si possono eseguire su di essa.
\end{itemize}
\end{definizione}

\subsubsection{Tipi in C}

\begin{itemize}
    \item \textbf{Tipi di base}: interi (standard ed estesi) e floating-point
          (reali e complessi).
    \item \textbf{Tipi enumerati}.
    \item Il tipo \lstinline|void|.
    \item \textbf{Tipi derivati}: puntatori, array, strutture, unioni, funzioni.
\end{itemize}

I tipi di base e i tipi enumerati formano i \emph{tipi aritmetici}; i tipi
aritmetici e i puntatori formano i \emph{tipi scalari}; array e strutture sono
detti \emph{tipi aggregati}. Un \emph{tipo funzione} descrive l'interfaccia di
una funzione: il tipo del valore restituito ed eventualmente i tipi dei
parametri.

\subsection{Interi}

I tipi interi con segno sono cinque (molti hanno sinonimi):
\begin{center}
\begin{tabular}{ll}
\toprule
\textbf{tipo} & \textbf{sinonimi} \\
\midrule
\lstinline|signed char| & \\
\lstinline|int| & \lstinline|signed|, \lstinline|signed int| \\
\lstinline|short| & \lstinline|short int|, \lstinline|signed short|, \lstinline|signed short int| \\
\lstinline|long| & \lstinline|long int|, \lstinline|signed long|, \lstinline|signed long int| \\
\lstinline|long long| (C99) & \lstinline|long long int|, \lstinline|signed long long|, \lstinline|signed long long int| \\
\bottomrule
\end{tabular}
\end{center}

C definisce solo le dimensioni minime: \lstinline|short| occupa almeno 2 byte,
\lstinline|long| almeno 4 e \lstinline|long long| almeno 8. I tipi possono essere
più grandi, ma devono rispettare l'ordine
\[
  \texttt{sizeof(short)} \leq \texttt{sizeof(int)} \leq \texttt{sizeof(long)}
  \leq \texttt{sizeof(long long)} .
\]

\subsubsection{Booleani}

In C non esistono \emph{true} e \emph{false}: $0$ è falso e qualsiasi valore
diverso da $0$ (ad esempio $1$, $-25$, $123456$) è vero.

\begin{codice}
if (3)                      if (0)
    printf("YES\n");            printf("YES\n");
else                        else
    printf("NO\n");             printf("NO\n");
\end{codice}

Il primo \lstinline|if| stampa \texttt{YES}, il secondo \texttt{NO}.

C99 ha introdotto il tipo intero senza segno \lstinline|_Bool|: \emph{vero} è
codificato come $1$, \emph{falso} come $0$. Includendo l'header
\lstinline|stdbool.h| si possono usare anche \lstinline|bool|, \lstinline|true| e
\lstinline|false|: \lstinline|bool| è sinonimo di \lstinline|_Bool| e
\lstinline|true|, \lstinline|false| sono costanti uguali a $1$ e $0$.

\subsubsection{Caratteri}

Anche \lstinline|char| è un tipo intero standard, ma il nome \lstinline|char| è
sinonimo di \lstinline|signed char| oppure di \lstinline|unsigned char| a seconda
del compilatore. Occupa 1 byte e i caratteri corrispondono ai codici della
tabella ASCII.

Con le variabili \lstinline|char| si può fare aritmetica: sta al programmatore
decidere se il numero contenuto è un codice di carattere o altro.

\begin{codice}
char ch = 'A';       // variabile di tipo char
printf("The character %c has the character code %d.\n", ch, ch);
printf("%c", ch + 1);
\end{codice}
\begin{risultato}
The character A has the character code 65.
B
\end{risultato}

\subsection{Sistemi di numerazione}

\begin{definizione}
Il sistema \emph{binario} (base 2) rappresenta i numeri con due soli simboli,
$0$ e $1$; ogni cifra è un \emph{bit}. Quasi tutti i calcolatori lo usano
internamente perché è semplice da realizzare con circuiti a porte logiche.
\end{definizione}

\begin{definizione}
Il sistema \emph{esadecimale} (base 16) usa sedici simboli: le cifre $0$--$9$ per
i valori da zero a nove e le lettere \texttt{A, B, C, D, E, F} (o minuscole) per
i valori da dieci a quindici. Ogni cifra esadecimale corrisponde a 4 bit (un
\emph{nibble}, metà di un byte da 8 bit): è una rappresentazione più
leggibile dei valori binari. Il sistema \emph{ottale} (base 8) usa le cifre
$0$--$7$.
\end{definizione}

\begin{center}
\begin{tabular}{rcc @{\qquad} rcc}
\toprule
\textbf{Bin.} & \textbf{Esa.} & \textbf{Dec.} &
\textbf{Bin.} & \textbf{Esa.} & \textbf{Dec.} \\
\midrule
0    & 0 & 0  & 1000 & 8 & 8 \\
1    & 1 & 1  & 1001 & 9 & 9 \\
10   & 2 & 2  & 1010 & A & 10 \\
11   & 3 & 3  & 1011 & B & 11 \\
100  & 4 & 4  & 1100 & C & 12 \\
101  & 5 & 5  & 1101 & D & 13 \\
110  & 6 & 6  & 1110 & E & 14 \\
111  & 7 & 7  & 1111 & F & 15 \\
\bottomrule
\end{tabular}
\end{center}

\subsubsection{Da binario/esadecimale a decimale}

Un numero con cifre $d_n \dots d_1 d_0$ in base $b$ vale
$d_n b^n + d_{n-1} b^{n-1} + \dots + d_1 b^1 + d_0 b^0$ (in base 2:
$d_n 2^n + \dots + d_1 2^1 + d_0 2^0 = N_{10}$).

\begin{esempio}
$1001_2 = 1 \cdot 2^3 + 0 \cdot 2^2 + 0 \cdot 2^1 + 1 \cdot 2^0 = 9$ \quad e \quad
$\texttt{002B}_{16} = 2 \cdot 16 + 11 = 43$.
\end{esempio}

\subsubsection{Da decimale a binario/esadecimale}

Si divide ripetutamente il numero per la base e si leggono i resti
\emph{dal basso verso l'alto}. Per l'ottale si usa lo stesso algoritmo con base 8.

\begin{esempio}
$156$ in binario (divisioni per 2): i quozienti successivi sono
$78, 39, 19, 9, 4, 2, 1, 0$ e i resti $0, 0, 1, 1, 1, 0, 0, 1$. Letti dal basso:
$156 = 10011100_2$.

$1565$ in esadecimale (divisioni per 16): $1565 = 97 \cdot 16 + 13$ (resto
$13 = \texttt{d}$), $97 = 6 \cdot 16 + 1$ (resto $1$), $6 = 0 \cdot 16 + 6$
(resto $6$). Quindi $1565 = \texttt{61d}_{16}$.
\end{esempio}

Con $N$ bit si rappresentano i numeri da $0$ a $2^N - 1$; con 8 bit (un byte)
da $0$ a $255$ (da \texttt{0000 0000} a \texttt{1111 1111}). Servono però anche i
numeri negativi, e per questo esistono rappresentazioni diverse: segno e modulo,
complemento a due.

\subsection{Rappresentazione degli interi in memoria}

\subsubsection{Segno e modulo}

Un bit (di solito il più a sinistra) indica il segno: $0$ per i positivi, $1$ per
i negativi; i bit restanti rappresentano il modulo. Ad esempio \texttt{1001 1000}
vale $-24$ (mentre se si usasse solo per numeri positivi varrebbe $152$).

Con $n$ bit si rappresentano i valori da $-(2^{n-1} - 1)$ a $2^{n-1} - 1$, e
$\pm 0$: con 8 bit da $-127$ a $+127$.

\begin{osservazione}
Problema: lo zero ha due rappresentazioni, \texttt{0000 0000} ($+0$) e
\texttt{1000 0000} ($-0$). Una soluzione è il complemento a due.
\end{osservazione}

\subsubsection{Complemento a due}

\begin{center}
\begin{tabular}{ccc}
\toprule
\textbf{Valore binario} & \textbf{Complemento a due} & \textbf{Senza segno} \\
\midrule
\texttt{00000000} & $0$    & $0$   \\
\texttt{00000001} & $1$    & $1$   \\
$\vdots$ & $\vdots$ & $\vdots$ \\
\texttt{01111110} & $126$  & $126$ \\
\texttt{01111111} & $127$  & $127$ \\
\texttt{10000000} & $-128$ & $128$ \\
\texttt{10000001} & $-127$ & $129$ \\
\texttt{10000010} & $-126$ & $130$ \\
$\vdots$ & $\vdots$ & $\vdots$ \\
\texttt{11111110} & $-2$   & $254$ \\
\texttt{11111111} & $-1$   & $255$ \\
\bottomrule
\end{tabular}
\end{center}

Nel complemento a due esiste un solo zero, \texttt{00000000}. Per cambiare segno a
un numero (positivo o negativo) si \textbf{invertono tutti i bit e poi si somma 1}.

\begin{esempio}
\begin{itemize}
    \item \texttt{0000 0101} (valore $5$); invertendo: \texttt{1111 1010};
          sommando $1$: \texttt{1111 1011}, che vale $-5$.
    \item Da \texttt{1111 1011}: invertendo \texttt{0000 0100}, sommando $1$
          \texttt{0000 0101}, cioè di nuovo $5$.
\end{itemize}
\end{esempio}

Con $n$ bit si rappresentano i valori da $-2^{n-1}$ a $2^{n-1} - 1$: non c'è
``$-0$'', quindi si rappresenta un numero negativo in più. Con 8 bit, da $-128$
(\texttt{1000 0000}) a $+127$ (\texttt{0111 1111}).

\begin{osservazione}
La regola ``inverti e somma 1'' non funziona per $-128$, perché $128$ non è
rappresentabile con 8 bit in complemento a due. D'ora in poi si usa il complemento
a due.
\end{osservazione}

\subsubsection{Endianness}

\begin{definizione}
L'\emph{endianness} indica l'ordine usato per interpretare numericamente una
sequenza di byte della memoria come una parola più grande (e l'ordine di
trasmissione dei byte su un collegamento digitale). Le parole possono essere in
formato \emph{big-endian} o \emph{little-endian}, a seconda che si parta
dall'estremità più significativa (big end) o meno significativa (little end).
\end{definizione}

Esempi: i mainframe IBM z/Architecture e i Motorola 68000 usano big-endian, i
processori Intel x86 little-endian.

La memoria è un array di byte; ogni byte ha il suo \emph{indirizzo logico}
(un numero positivo), cioè l'indirizzo al quale un elemento appare risiedere dal
punto di vista di un programma in esecuzione. Per architetture a 64 bit il limite
superiore è $2^{64} - 1$.

\begin{esempio}
Il numero $2$ scritto su 4 byte, a partire dall'indirizzo $1500$:
\begin{center}
\begin{tabular}{lcccc}
\toprule
 & 1500 & 1501 & 1502 & 1503 \\
\midrule
big-endian    & \texttt{00} & \texttt{00} & \texttt{00} & \texttt{02} \\
little-endian & \texttt{02} & \texttt{00} & \texttt{00} & \texttt{00} \\
\bottomrule
\end{tabular}
\end{center}
\end{esempio}

\begin{esempio}
Sulla stringa di 4 bit \texttt{0010}, in complemento a due: in big-endian vale $2$;
in little-endian si leggono i bit in ordine inverso (\texttt{0100}) e vale $4$.
Sulla stringa \texttt{1010}: in big-endian vale $-6$, in little-endian
(\texttt{0101}) vale $5$.
\end{esempio}

\subsection{Dimensione e limiti dei tipi}

L'operatore \lstinline|sizeof| restituisce la dimensione esatta in byte di un tipo
o di una variabile: \lstinline|sizeof(tipo)| oppure \lstinline|sizeof espressione|
(in questo caso conta il tipo dell'espressione).

\begin{codice}
int iIndex;
iIndex = 1000;
// sizeof(int) e sizeof(iIndex) valgono entrambi 4
\end{codice}

I limiti dei tipi interi per il proprio compilatore sono nell'header
\lstinline|limits.h|, che definisce macro come \lstinline|INT_MIN|,
\lstinline|INT_MAX|, \lstinline|UINT_MAX|, \lstinline|CHAR_MIN|,
\lstinline|CHAR_MAX|.

\begin{codice}
#include <stdio.h>
#include <limits.h>

int main() {
    printf(" char %zu %d %d\n", sizeof(char), CHAR_MIN, CHAR_MAX);
    printf(" int %zu %d %d\n", sizeof(int), INT_MIN, INT_MAX);
    return 0;
}
\end{codice}
\begin{risultato}
 char 1 -128 127
 int 4 -2147483648 2147483647
\end{risultato}

\begin{osservazione}
Il risultato di \lstinline|sizeof| ha tipo \lstinline|size_t|: per stamparlo è
corretto \lstinline|%zu| (nelle slide compare \lstinline|%d|). I valori stampati
dipendono dalla piattaforma (qui un caso tipico).
\end{osservazione}

\subsection{Numeri in virgola mobile}

Per rappresentare i non interi, con la virgola in qualsiasi posizione, C ha tre
tipi floating-point standard:
\begin{itemize}
    \item \lstinline|float|: precisione singola;
    \item \lstinline|double|: precisione doppia;
    \item \lstinline|long double|: precisione estesa.
\end{itemize}

\subsubsection{Precisione}

Un valore in virgola mobile si memorizza solo con una \emph{precisione} limitata,
determinata dal formato binario e dalla memoria usata. La precisione si esprime in
cifre significative: ad esempio con sei cifre significative la conversione in un
numero decimale di sei cifre restituisce le sei cifre originali. La posizione della
virgola e gli zeri iniziali o finali non contano: $123\,456\,000$ e
$0{,}00123456$ si possono entrambi memorizzare in un tipo con precisione a sei
cifre.

In C le operazioni aritmetiche tra numeri in virgola mobile sono eseguite
internamente con precisione double o superiore.

\begin{codice}
float height = 1.2345, width = 2.3456;   // precisione singola
double area = height * width;            // il calcolo è fatto in double
\end{codice}

Se il risultato è assegnato a una variabile \lstinline|float|, viene arrotondato
se necessario. L'header \lstinline|float.h| definisce macro come
\lstinline|FLT_MIN|, \lstinline|FLT_MAX| e \lstinline|FLT_DIG|, che indicano
intervallo di valori e precisione del tipo \lstinline|float|.

\subsubsection{Notazione E e formato IEEE 754}

La \emph{notazione E} è la rappresentazione testuale della notazione scientifica:
\texttt{1.234e+56} significa $1{,}234 \times 10^{56}$.

Nel formato \emph{IEEE 754} ogni numero finito è descritto da tre interi: un segno
$s$ (zero o uno), un \emph{significando} (o mantissa) $c$ e un esponente $q$. Il
valore è
\[
  (-1)^s \times c \times b^q ,
\]
dove $b$ è la base (2 o 10). Ad esempio con base 10, segno $1$ (negativo),
significando $12345$ ed esponente $-3$ il valore è $-12{,}345$.

Il formato 754 a \emph{precisione singola} usa: 1 bit di segno, 8 bit di esponente,
23 bit di significando. Il valore è
\[
  (-1)^{\text{segno}} \times \Bigl( 1 + \sum_{i=1}^{23} b_{23-i}\, 2^{-i} \Bigr)
  \times 2^{(e - 127)} .
\]

\begin{esempio}
Con segno $0$, esponente \texttt{01111100} ($e = 124$) e frazione
\texttt{010\dots0} (cioè $b_{21} = 1$, tutti gli altri bit nulli) si ha
$1 + 2^{-2} = 1{,}25$ e $2^{124 - 127} = 2^{-3}$, quindi
\[
  \text{valore} = (+1) \times 1{,}25 \times 2^{-3} = +0{,}15625 .
\]
\end{esempio}

Con la \emph{precisione doppia} (\lstinline|double|) si usano 1 bit di segno,
11 di esponente e 52 di frazione; la \emph{precisione estesa}
(\lstinline|long double|, 80 bit) ha 1 bit di segno, 15 di esponente, 1 bit di
parte intera e 63 di frazione.

\subsubsection{Errori di arrotondamento}

L'errore di arrotondamento è intrinseco nel calcolo in virgola mobile.

\begin{codice}
#include <stdio.h>

int main() {
    int a = 16777217;
    float b = a;
    printf("%f\n", b);
}
\end{codice}
\begin{risultato}
16777216.000000
\end{risultato}

Il valore $16777217 = 2^{24} + 1$ non è rappresentabile esattamente con un
\lstinline|float|. Allo stesso modo conviene tenere in \lstinline|double| anche
il risultato di un prodotto come \lstinline|height * width|: salvarlo in un
\lstinline|float| fa perdere precisione.

\begin{osservazione}
Il modo più semplice per evitare di accumulare errori è usare numeri in virgola
mobile ad alta precisione, cioè \lstinline|double| al posto di \lstinline|float|.
Sulle CPU moderne non c'è (quasi) penalizzazione di tempo, ma un \lstinline|double|
occupa il doppio della memoria.
\end{osservazione}

\subsection{Tipi enumerati}

Le \emph{enumerazioni} sono tipi interi definiti dal programmatore. La definizione
inizia con la parola chiave \lstinline|enum|, eventualmente seguita da un
identificatore, e contiene l'elenco dei possibili valori con un nome per ciascuno:
\begin{codice}
enum [identificatore] { lista-enumeratori };
\end{codice}

\begin{esempio}
\begin{codice}
enum color { black, red, green, yellow, blue, white = 7, gray };
\end{codice}
Le costanti elencate valgono $0, 1, 2, 3, 4, 7, 8$ (se non si assegna un valore,
si continua dal precedente più uno).
\end{esempio}

\begin{codice}
enum color fgColor = blue,          // due variabili
           bgColor = yellow;        // di tipo enum color

void setFgColor(enum color fgc);    // funzione con parametro enum color
\end{codice}

Con le variabili di tipo enumerato si può fare aritmetica ordinaria
(\lstinline|red + red| vale $2$). Costanti diverse della stessa enumerazione
possono avere lo stesso valore:
\begin{codice}
enum signals { OFF, ON, STOP = 0, GO = 1, CLOSED = 0, OPEN = 1 };
\end{codice}

\begin{esempio}
\begin{codice}
#include <stdio.h>

enum week { sunday, monday, tuesday, wednesday, thursday, friday, saturday };

int main() {
    enum week today;
    today = wednesday;
    printf("Day %d", today + 1);
    return 0;
}
\end{codice}
\begin{risultato}
Day 4
\end{risultato}
\end{esempio}

\begin{osservazione}
Conviene usare sempre un'enumerazione quando una variabile (in particolare un
parametro di funzione) può assumere solo uno tra pochi valori possibili: si evitano
errori dovuti a costanti non valide, si documentano i valori ammessi ed è più
mnemonico degli interi.
\end{osservazione}

\subsection{Il tipo void}

Il tipo \lstinline|void| indica che non è disponibile alcun valore. Non si possono
dichiarare variabili o costanti di questo tipo, ma si usa:
\begin{itemize}
    \item nelle dichiarazioni di funzioni;
    \item nelle espressioni \lstinline|void|;
    \item nei puntatori a \lstinline|void| (si vedranno con i puntatori).
\end{itemize}

Una funzione che non restituisce nulla ha tipo \lstinline|void|:
\lstinline|void error(int a) {}|. \lstinline|void| nella lista dei parametri di un
prototipo indica che la funzione non ha parametri:
\lstinline|void printMenu(void) {}|. Il compilatore segnala un errore se si
chiama \lstinline|printMenu(3)|.

\end{document}
